\documentclass[14pt,a4paper,oneside]{extbook}
\usepackage[utf8]{vietnam}
\usepackage{amsmath,amssymb,mathrsfs,fancyhdr,enumerate,multirow,makecell}
\usepackage{tkz-euclide,tikz-3dplot,tikz,tkz-tab}
\usetikzlibrary{arrows.meta,calc,intersections,angles,quotes,backgrounds,shapes.geometric,patterns,through,shadows,positioning}
\usepackage{yhmath}
\usepackage[top=0.7cm, bottom=0.7cm, left=1.0cm, right=1.0cm, headheight=16pt, headsep=3mm, footskip=5mm]{geometry}
\usepackage[hidelinks,unicode]{hyperref}
\usepackage{qrcode}
\usepackage[many]{tcolorbox}
\usepackage{fontawesome}

\pagestyle{fancy}
\fancyhf{}
\fancyhead[L]{\footnotesize\bfseries\color{teal!85!black}\faBook~SỔ TAY KIẾN THỨC TOÁN 5}
\fancyhead[R]{\footnotesize\bfseries\color{blue!85!black}Lớp Toán Thầy Lĩnh -- loptoanthaylinh.vn}
\fancyfoot[C]{\footnotesize\bfseries\color{gray!80!black}-- Sổ tay kiến thức: Bài 48. Luyện tập chung đơn vị đo thể tích --}
\renewcommand{\headrulewidth}{0.4pt}

\begin{document}

% --- BANNER TIÊU ĐỀ BÀI HỌC ---
\noindent
\begin{tcolorbox}[
  enhanced,
  colback=teal!8!white,
  colframe=teal!80!black,
  arc=3mm,
  boxrule=1.1pt,
  width=\linewidth,
  halign=center,
  top=2.2mm, bottom=2.2mm,
  left=3mm, right=3mm,
  drop fuzzy shadow
]
  {\large\bfseries\color{teal!95!black}\faBookmark\ BÀI 48. LUYỆN TẬP CHUNG ĐƠN VỊ ĐO THỂ TÍCH (CHỦ ĐỀ 8)}\\[2pt]
  {\scriptsize\color{black!75}\itshape Kỹ năng thành thạo đổi đơn vị thể tích sang số thập phân, so sánh và giải toán bể nước trên 1 trang A4}
\end{tcolorbox}

\vspace{3.2mm}
% --- KHỐI 1: BẢNG SO SÁNH BA HỆ ĐƠN VỊ ĐO LƯỜNG ---
\noindent
\begin{tcolorbox}[
  enhanced,
  title={\footnotesize\bfseries\color{white}\faTable\ 1. Bảng so sánh bước nhảy giữa 3 hệ đơn vị đo lường},
  colback=blue!2!white,
  colframe=blue!75!black,
  arc=2.5mm, boxrule=0.8pt, top=1.8mm, bottom=1.8mm, left=2mm, right=2mm
]
  \scriptsize
  \centering
  \begin{tabular}{|l|c|c|l|}
    \hline
    \textbf{Hệ đo lường} & \textbf{Hai đơn vị liền kề} & \textbf{Số chữ số mỗi cột} & \textbf{Ví dụ minh họa} \\ \hline
    \textbf{Độ dài} & Gấp/kém \textbf{10 lần} & \textbf{1 chữ số} & $3\text{ m } 5\text{ dm} = 3,5\text{ m}$ \\ \hline
    \textbf{Diện tích} & Gấp/kém \textbf{100 lần} & \textbf{2 chữ số} & $3\text{ m}^2\ 5\text{ dm}^2 = 3,05\text{ m}^2$ \\ \hline
    \textbf{Thể tích} & Gấp/kém \textbf{1 000 lần} & \textbf{3 chữ số} & $3\text{ m}^3\ 5\text{ dm}^3 = 3,005\text{ m}^3$ \\ \hline
  \end{tabular}
\end{tcolorbox}

\vspace{3mm}

% --- KHỐI 2: KỸ NĂNG SO SÁNH VÀ TÍNH TOÁN VỚI SỐ ĐO THỂ TÍCH ---
\noindent
\begin{minipage}[t]{0.52\textwidth}
  \begin{tcolorbox}[
    enhanced,
    title={\footnotesize\bfseries\color{white}\faBalanceScale\ 2. So sánh các số đo thể tích},
    colback=teal!2!white,
    colframe=teal!75!black,
    arc=2.5mm, boxrule=0.8pt, top=1.8mm, bottom=1.8mm, left=2.2mm, right=2.2mm
  ]
    \footnotesize
    So sánh $2\text{ m}^3\ 15\text{ dm}^3$ và $2,15\text{ m}^3$:\\
    $\triangleright$ Cách 1: Đổi ra $\text{dm}^3$:\\
    \quad $2\text{ m}^3\ 15\text{ dm}^3 = 2\,015\text{ dm}^3$.\\
    \quad $2,15\text{ m}^3 = 2\,150\text{ dm}^3$.\\
    Vì $2\,015 < 2\,150$ nên $2\text{ m}^3\ 15\text{ dm}^3 < 2,15\text{ m}^3$.\\
    $\triangleright$ Cách 2: Đổi ra $\text{m}^3$: $2,015\text{ m}^3 < 2,15\text{ m}^3$.
  \end{tcolorbox}
\end{minipage}%
\hfill
\begin{minipage}[t]{0.46\textwidth}
  \begin{tcolorbox}[
    enhanced,
    title={\footnotesize\bfseries\color{white}\faCalculator\ 3. Phép tính với số đo thể tích},
    colback=purple!2!white,
    colframe=purple!70!black,
    arc=2.5mm, boxrule=0.8pt, top=1.8mm, bottom=1.8mm, left=2mm, right=2mm
  ]
    \footnotesize
    $\bullet$ Cộng, trừ, nhân, chia như các số thập phân thông thường kèm tên đơn vị.\\
    \textit{VD:} $25,4\text{ m}^3 + 12,8\text{ m}^3 = 38,2\text{ m}^3$.\\
    $4,5\text{ dm}^3 \times 4 = 18\text{ dm}^3$.\\
    $12\text{ cm}^3 : 5 = 2,4\text{ cm}^3$.
  \end{tcolorbox}
\end{minipage}

\vspace{3mm}

% --- KHỐI 3: SƠ ĐỒ TƯ DUY ---
\noindent
\begin{tcolorbox}[
  enhanced,
  title={\footnotesize\bfseries\color{white}\faSitemap\ Sơ đồ tư duy: Kỹ thuật đổi số đo thể tích},
  colback=yellow!2!white,
  colframe=teal!85!black,
  arc=3mm, boxrule=0.9pt, top=1.8mm, bottom=1.8mm, left=1.5mm, right=1.5mm,
  halign=center
]
  \begin{tikzpicture}[
    scale=0.92,
    every node/.style={transform shape},
    b1/.style={rectangle, rounded corners=4pt, fill=blue!15, draw=blue!70!black, text=blue!90!black, font=\bfseries\scriptsize, inner sep=4.5pt, align=center},
    b2/.style={rectangle, rounded corners=4pt, fill=orange!15, draw=orange!70!black, text=orange!90!black, font=\bfseries\scriptsize, inner sep=4.5pt, align=center}
  ]
    \node[b1] (n1) at (-4.8, 0) {\textbf{ĐƠN VỊ LỚN $\rightarrow$ BÉ: Nhân $1\,000$}\\[2pt]Dịch dấu phẩy sang \textbf{phải} 3 hàng};
    \node[b2] (n2) at (4.8, 0) {\textbf{ĐƠN VỊ BÉ $\rightarrow$ LỚN: Chia $1\,000$}\\[2pt]Dịch dấu phẩy sang \textbf{trái} 3 hàng};

    \draw[teal!85!black, very thick, <->] (n1) -- (n2);
  \end{tikzpicture}
\end{tcolorbox}

\vspace{3mm}

% --- KHỐI 4: QUY TẮC BỂ NƯỚC & BẪY TRẮC NGHIỆM ---
\noindent
\begin{minipage}[t]{0.54\textwidth}
  \begin{tcolorbox}[
    enhanced,
    title={\footnotesize\bfseries\color{white}\faTint\ Cầu nối thần kỳ: $\text{dm}^3$ và Lít},
    colback=green!3!white,
    colframe=green!60!black,
    arc=2.5mm, boxrule=0.8pt, top=1.8mm, bottom=1.8mm, left=2.2mm, right=2.2mm
  ]
    \footnotesize
    Trong các bài toán bể chứa nước:\\
    $\bullet$ Đổi các kích thước (dài, rộng, cao) ra đơn vị \textbf{$\text{dm}$}.\\
    $\bullet$ Khi đó thể tích tính ra theo $\text{dm}^3$ chính là \textbf{số lít nước} luôn ($1\text{ dm}^3 = 1\text{ lít}$)!\\
    Không cần mất công đổi lằng nhằng ở bước cuối.
  \end{tcolorbox}
\end{minipage}%
\hfill
\begin{minipage}[t]{0.44\textwidth}
  \begin{tcolorbox}[
    enhanced,
    title={\footnotesize\bfseries\color{white}\faExclamationTriangle\ Mẹo tránh bẫy trắc nghiệm},
    colback=red!2!white,
    colframe=red!70!black,
    arc=2.5mm, boxrule=0.8pt, top=1.8mm, bottom=1.8mm, left=2.2mm, right=2.2mm
  ]
    \footnotesize
    $\bullet$ Viết thiếu chữ số $0$ khi đổi sang $\text{m}^3$: Hàng $\text{dm}^3$ phải đủ $3$ chữ số thập phân.\\
    $\bullet$ Khác nhau giữa số mũ $2$ (diện tích) và số mũ $3$ (thể tích).
  \end{tcolorbox}
\end{minipage}

\end{document}
